\cotefile{148}
\scanpage{44}
\mspage{1}
% archive sheet number: 43
\uline{Groupoïdes de Teichmüller}.

\textcircled{1} Cas topologique.

% slanted
\margin{Finalement, tant qu'on ne travaille pas au point de vue arithmétique, il semble que c'est le gros $ {\unsure{\mathscr{T}|\mathcal{T}}}^{\pm}$ qui sera le plus utilisé, \deleted{\unsure{tandis}} c'est lui qui \unsure{intervient|apparaît} dans \unsure{ici} la factorisation la plus simple $\mathscr{T}$, et le \unsure{gpde} orienté \unsure{des surfaces|s'appelle} $\mathscr{T}^{+}$.}
La variante la plus simple est
\[
\mathscr{T}_{\mathrm{top}}^{+}\quad\text{ou}\quad \mathscr{T}_{\mathrm{top}}
\]
(\uline{groupoïde de Teichmüller} \uuline{orienté} \uline{topologique}),

\uline{objets}: surfaces \uline{à bord} compactes orientées

\uline{flèches}: classes d'isotopie d'iso (conservant l'orientation)

Variante immédiate
\[
\mathscr{T}_{\mathrm{top}}^{\pm}
\]
(\uline{groupoïde de Teichmüller} \uuline{biorienté} \uline{topologique}) :
mêmes objets, flèches sont les classes d'isotopie d'homéomorphismes (ne conservant pas néc. l'orientation).

Soit $[\pm1]$ le groupoïde ayant un seul objet $e$ (par exemple $e=\unsure{\varnothing|\phi}$) et comme groupe d'autom $\pm1$. On a un morphisme de signature
\diagram{
\begin{tikzcd}[ampersand replacement=\&]
\mathscr{T}_{\mathrm{top}}^{\pm}\arrow[r,"\unsure{\chi|X}"] \& [\pm1]
\end{tikzcd}
}
% slanted
\margin{Ça ne \unsure{vaut|va}, \unsure{si on se borne|sinon se borner} à des surfaces connexes !}
dont la \uline{fibre} \deleted{\unsure{est}} (ou noyau) est $\mathscr{T}_{\mathrm{top}}^{+}=\mathscr{T}_{\mathrm{top}}$.

La variante la plus complexe est
\[
S_{*}\mathscr{T}{\unsure{D|\mathbb D}}_{\mathrm{top}}^{+}\quad\text{ou}\quad S_{*}\mathscr{T}D_{\mathrm{top}},
\]
``\uline{groupoïde de Teichmüller} \inserted{\uline{orienté},} \uline{spécial dévissé topologique}''
